\paragraph{A finite second-descent bound from the six isogeny classes.}
Write \(A=E_b\), \(B=E_{-4b}\), and let
\(\phi:A\to B\), \(\psi:B\to A\) be the actual dual maps
of A.150.16.  In particular \(\psi\phi=[2]_A\) and
\(\phi\psi=[2]_B\).  We prove
\[
                      |\operatorname{Sel}^{2}(A/K)|\leq8
\tag{A.154.S1}
\]
using only the two already proved groups in A.150.6
and the explicit finite calculation below.

For \(d\in\mathcal S_b\), put \(B_d=b/d\) and
\[
 C_d:\ N^2=dU^4+B_dV^4,\qquad
 \rho_d:C_d\longrightarrow A,\quad
 (U:V:N)\longmapsto
 \left(\frac{dU^2}{V^2},\frac{dUN}{V^3}\right).
\tag{A.154.S2}
\]
As proved by the actual inverse formula in A.150.18--20,
over \(K(\sqrt d)\) this is the map \(\psi:B\to A\).
Changing the sign of \(\sqrt d\) translates the source
by \(T'=(0,0)\in B[\psi]\).  Thus
\[
 \mathcal S_b=\operatorname{Sel}^{\psi}(B/K)
     \longrightarrow \Sha(B/K)[\psi],\qquad
 d\longmapsto [C_d],
\tag{A.154.S3}
\]
is induced by the literal inclusion
\(B[\psi]=\{O,T'\}\hookrightarrow B\):
the cocycle of \(C_d\) is \(g\mapsto\kappa_d(g)T'\).
This identifies its Jacobian with \(B\), including the
specified kernel point.  There is no unspecified
identification of genus-one models here.

\emph{The tangent function and its divisor class.}
Choose \(n_d,z_d,v_d\in K\) with
\[
                  n_d^2=dz_d^2+B_dv_d^2 .
\tag{A.154.S4}
\]
The six choices used below have \(v_d=1\), and all
three coordinates are nonzero.  The map
\[
 C_d\longrightarrow Q_d,\qquad
  (U:V:N)\longmapsto(N:U^2:V^2),\qquad
 Q_d:\ n^2=dz^2+B_dv^2
\tag{A.154.S5}
\]
is the degree-two quotient by
\(j(U:V:N)=(-U:V:N)\).
The line tangent to \(Q_d\) at \((n_d:z_d:v_d)\)
pulls back to
\[
 \ell_d=n_dN-dz_dU^2-B_dv_dV^2.
\]
The exact polynomial identity is
\[
 \ell_d\bigl(n_dN+dz_dU^2+B_dv_dV^2\bigr)
          =b(v_dU^2-z_dV^2)^2.
\tag{A.154.S6}
\]
Its geometric divisor is twice the degree-two fibre
of (A.154.S5).  Consequently
\[
 f_d=\ell_d/V^2,\qquad
       \operatorname{div}(f_d)=2D_d,\qquad
       D_d=F_j-F_\iota\in\operatorname{Div}^0(C_d),
\tag{A.154.S7}
\]
where \(F_j\) is that conic fibre and \(F_\iota\)
is the fibre at infinity of the hyperelliptic map
\((U:V:N)\mapsto(U:V)\).
Its involution is \(\iota(U:V:N)=(U:V:-N)\).
For a degree-two map from a genus-one curve the
involution is \(P\mapsto F-P\) in the corresponding
degree-two divisor class.  Thus \(j\iota\) is
translation by \([F_j-F_\iota]\).
The literal formula
\[
                    j\iota(U:V:N)=(-U:V:-N)
\]
fixes both functions in (A.154.S2), and is its
nonidentity deck transformation.  Under its already
specified Jacobian \(B\), it is translation by \(T'\).
Therefore \([D_d]=T'\in\operatorname{Pic}^0(C_d)(K)\).
This is also the divisor calculation of
\cite[Lemma 6.1]{AgentFisherBinaryQuartics2022}.

For \(e\in\mathcal S_b\), the quaternion
\((e,f_d)\) is unramified on the proper curve,
because (A.154.S7) has even divisor.
Over \(K(\sqrt e)\), the norm of the rational
divisor \(D_d\) is \(2D_d=\operatorname{div}(f_d)\).
The connecting map from divisors to the Brauer group
therefore sends the cocycle \(\kappa_e T'\) to
\((e,f_d)\).  This proves, with its actual cocycle
and function, the formula
\[
 (-1)^{\,2\langle[C_d],[C_e]\rangle_{\rm CT}}
       =\prod_w(e,f_d(P_w))_w,\qquad P_w\in C_d(K_w).
\tag{A.154.S8}
\]
Here the exponent means the identification
\(\tfrac12\Z/\Z\simeq\{\pm1\}\).
Equivalently the additive invariant is the sum
of the quaternion invariants.
This is the homogeneous-space definition and
cyclic-algebra map of
\cite[Section 4, equation (10), and Lemma 5.1]
{AgentFisherBinaryQuartics2022}.
At a chart with \(V=0\), use \(\ell_d/U^2\);
the two functions differ by a square wherever
both charts are defined, so this is the same
Brauer evaluation.  This accounts for every
chart used in the finite certificate.

\emph{Six literal conic points and the finite calculation.}
Let \(d_1,\ldots,d_6\) be the ordered basis of
\(\mathcal S_b\) in A.150.6, let \(t=\sqrt2\),
and retain \(M\) from A.150.1.  The following table
gives \(n_{d_i}=M\nu_i\), \(z_{d_i}=M\zeta_i\),
\(v_{d_i}=1\).  Direct substitution verifies
\(\nu_i^2=d_i\zeta_i^2+1146(1+t)^2/d_i\).
\[
 \begin{array}{c|c|c}
 i&\nu_i&\zeta_i\\ \hline
 1&-50-26t&-18-17t\\
 2&-26(1+t)/7&(2+t)/7\\
 3&(90+37t)/7&(-262-78t)/161\\
 4&(182-163t)/21&(10596-8182t)/67137\\
 5&(-47-29t)/17&(4098363+2550620t)/133181111\\
 6&(6196+4761t)/3361&
           (584427298+410557178t)/26330689063
 \end{array}
\tag{A.154.S9}
\]
These points were found by a conic search; their
displayed equations, rather than an output about
class groups or Selmer groups, are the entire
input from that search.

For each row, multiply the coefficients of
\(\ell_d\) by a common denominator and remove
their greatest common divisor in
\(\mathcal O_K=\Z[t]\).
The Euclidean algorithm of A.150 proves that
this is legitimate.  The certificate retains
explicit integral Bezout coefficients showing
that the three resulting coefficients generate
the unit ideal.  Multiplying \(f_d\) by this
single global constant does not change the
product in (A.154.S8), by global quaternion
reciprocity.

The only rational finite primes used are
\[
             2,3,5,7,13,23,139,191,1119169.
\tag{A.154.S10}
\]
Every \(d_i\) and \(b/d_i\) is integral and
supported on this list; this follows either
from A.150.22--23 or literal division in
\(\Z[t]\).  The second symbol entries in
(A.154.S8) use these \(d_i\), so they are units
outside this list.
At every other prime the quartic has smooth
genus-one reduction.  The primitive weighted
quadratic tangent cannot vanish identically
on that curve: its reduction has nonzero
coefficients and weighted degree two, whereas
the defining equation has weighted degree
four.  Its zero divisor has degree four.
The residue field has size \(q\geq11\), and
\[
                   q+1-2\sqrt q>4.
\]
There is consequently a residue point at
which the tangent is a unit.  One of \(U,V\)
is a unit; divide by its square and lift the
smooth point.  The resulting Hilbert symbol
of two units at an odd place is one.
This proves the vanishing at every omitted
finite place.  All \(d_i\) are totally positive,
so the two real symbols are also one.

Here are the exact local operations in the
retained certificate.  At a split odd prime
the residue map is \(t\mapsto t_w\), with its
unique Hensel lift from \(t_w^2=2\); at an inert
prime it is the actual field
\(\F_p[t]/(t^2-2)\).
For \(a=\varpi^\alpha a_0\), \(c=\varpi^\gamma c_0\)
over an odd residue field \(\F_q\), it uses
\[
 (a,c)_w=(-1)^{\alpha\gamma(q-1)/2}
              \overline a_0^{\,\gamma(q-1)/2}
              \overline c_0^{\,\alpha(q-1)/2}.
\tag{A.154.S11}
\]
This is the elementary quadratic norm test:
the unramified quadratic norm group consists
of the elements of even valuation; the
ramified quadratic norm on units is the
residue-square test.  Bilinearity and
\((a,-a)_w=1\) give the displayed formula.
For an inert place, the quadratic character
of a residue \(a+ct\) equals the Legendre
symbol of \(a^2-2c^2\) in \(\F_p\), since
this is the residue norm.

At the dyadic place use the squareclass basis
\((-1,5,\epsilon,\pi)\), where
\(\epsilon=1+\pi\) and \(\pi=\sqrt2\).
In additive \(\F_2\)-notation its Hilbert matrix is
\[
 \begin{pmatrix}
 0&0&1&1\\
 0&0&0&1\\
 1&0&1&1\\
 1&1&1&1
 \end{pmatrix}.
\tag{A.154.S12}
\]
For clarity, this matrix has a direct norm proof.
Two rational entries pair trivially after the
degree-two extension from \(\Q_2\).
Corestriction and
\(N_{K_2/\Q_2}(\epsilon)=-1\),
\(N_{K_2/\Q_2}(\pi)=-2\) give
\((\epsilon,-1)=(\pi,-1)=-1\) and
\((\epsilon,5)=1\), \((\pi,5)=-1\).
Here \(5\) defines the unramified quadratic
extension of \(\Q_2\), whose norms have even
valuation.  The remaining rational test
\((-1,-1)_{\Q_2}=-1\) follows since a sum of
two squares cannot be minus one in \(\Q_2\);
\((2,-1)_{\Q_2}=1\) follows from
\(2=N_{\Q_2(i)/\Q_2}(1+i)\).
The Steinberg relation
\((\epsilon,1-\epsilon)=(\epsilon,-\pi)=1\)
gives \((\epsilon,\pi)=-1\), and
\((a,a)=(a,-1)\) gives the two diagonal
entries.  These are every entry of (A.154.S12).
The basis is exhaustive: its sixteen products
give all sixteen squareclasses enumerated
modulo \(\pi^5\) in A.150.25.

All chosen local coordinates are retained
as exact rational pairs in \(K\).
For an ordinary odd square-root lift,
an approximation whose squared error has
valuation at least \(L+v(f)\) differs from
the actual root by valuation at least
\(L+v(f)/2\).
For an odd fourth-root lift the derivative
is a unit.  At two the certificate divides
out an even valuation, starts from a unit
square root modulo \(\pi^5\), and performs
Newton steps; its error bound on the root
is \(L+v_\pi(f)/2-2\).
Every tangent evaluation is checked stable
by a margin at least one at odd places and
five at two.  Thus the recorded signs are
symbols at actual local points, rather than
at unlifted residue points.

The resulting additive matrix on the
ordered basis \(d_1,\ldots,d_6\) is
\[
 \mathcal M=
 \begin{pmatrix}
 0&0&0&1&1&1\\
 0&0&0&1&1&1\\
 0&0&0&1&0&0\\
 1&1&1&0&0&1\\
 1&1&0&0&0&0\\
 1&1&0&1&0&0
 \end{pmatrix},
 \qquad \operatorname{rank}_{\F_2}\mathcal M=4 .
\tag{A.154.S13}
\]
Its radical is explicitly
\(\langle d_1d_2,d_3d_5d_6\rangle
 =\langle573,23\cdot1119169\rangle\):
the second product is literally
\(23\cdot7^2\cdot1119169\).
The literal row operations, six conic equations,
primitive Bezout identities, and every local
witness and sign are executed in
\path{code/harmonic_prym_isogeny_pairing_certificate.py}
and its retained output and data.

\emph{The exact Selmer maps giving the bound.}
The finite group schemes fit into
\[
 0\longrightarrow A[\phi]\longrightarrow A[2]
       \xrightarrow{\phi}B[\psi]\longrightarrow0.
\tag{A.154.S14}
\]
Its cohomology maps, with the local Kummer
conditions, give
\[
 \operatorname{Sel}^{\phi}(A/K)
       \longrightarrow\operatorname{Sel}^{2}(A/K)
       \xrightarrow{\phi_*}\operatorname{Sel}^{\psi}(B/K).
\tag{A.154.S15}
\]
The kernel of the first map is the boundary
of \(B[\psi](K)/\phi A[2](K)\).
Both rational kernels have order two,
\(\phi A[2](K)=0\), and this boundary is
the nonzero class \([-4b]\).
Since A.150 gives
\(|\operatorname{Sel}^{\phi}(A/K)|=4\),
its image has order two.
This is the whole kernel of the second
map.  To check the local condition in
that exactness assertion, take a global
lift in \(H^1(K,A[\phi])\) of a Selmer
class whose second image is zero.
At \(w\), write the original class
as \(\delta_2(P_w)\).
Its zero \(\psi\)-Kummer image says
\(P_w=\psi Q_w\) for \(Q_w\in B(K_w)\).
Subtracting \(\delta_\phi(Q_w)\) leaves
the boundary of \(B[\psi](K_w)\), which
is itself a local \(\phi\)-Kummer class.
Thus the global lift satisfies every
required local condition.

For \(\xi\in\operatorname{Sel}^{2}(A/K)\),
let \(\bar\xi\in\Sha(A/K)[2]\) be its
image.  The torsor image of \(\phi_*\xi\)
under (A.154.S3) is literally
\(\phi_*\bar\xi\in\Sha(B/K)\).
For every \([C_e]\) from (A.154.S3),
\(\psi_*[C_e]=0\).  The adjointness of
the actual dual isogenies therefore gives
\[
 \langle\phi_*\bar\xi,[C_e]\rangle_B
       =\langle\bar\xi,\psi_*[C_e]\rangle_A=0.
\tag{A.154.S16}
\]
The adjointness statement and its proof
by pulling back the divisor functions are
\cite[Theorem 1.2 and Section 2, pp.\ 187--188]
{AgentCasselsIsogenies1965}.
Thus \(\phi_*\operatorname{Sel}^{2}(A/K)\)
lies in the radical of the proved matrix
(A.154.S13), which has dimension two.
Combining this with the kernel of order
two in (A.154.S15) proves (A.154.S1).
No order of a computed full Selmer group,
and no finiteness hypothesis on a
Tate--Shafarevich group, enters this bound.
